\section{Stage 1 results: which guidance matters}
\label{sec:stage1}
\label{sec:stage1results}

All Stage-1 objective deltas are contrasts on the POC via a mixed-effects model
$y_{cps}=\mu+\tau_c+u_p+v_{cp}+\epsilon_{cps}$ with fixed cell effect $\tau_c$, random prompt
intercept $u_p$, and random cell-by-prompt slope $v_{cp}$; this is what separates the study from
per-generation $t$-tests, because five seeds on one prompt are not five independent points---with a
modest intraclass correlation the effective $n$ is roughly halved (Appendix~\ref{app:theory}).
Preference deltas are style-controlled Bradley--Terry utilities (Eq.~\ref{eq:bt}). Confirmatory
contrasts are Holm-corrected over the 12-test family, separately on each signal; the per-component
estimates are collected in Table~\ref{tab:s1effects} and visualized in Figure~\ref{fig:F1}.
For each research question we give the completed interpretation the paper will carry under
\emph{every} preregistered outcome; on real data, exactly one branch per question is kept.

\subsection{Does the skill work, and is it content?}
The headline contrast is \FULL{} vs.\ \NEUTRAL. Because \NEUTRAL{} is length-matched, a gain is
design \emph{content}, not prompt mass. We measure $\numslot{HSKILL-POC}$ standardized $\delta$ on
the POC and a \FULL{} win-rate of $\numslot{HSKILL-WIN}$ with BT utility $\numslot{HSKILL-BT}$; the
prompt-mass control \NEUTRAL{} vs.\ \NOSYS{} moves the POC by $\numslot{NEU-NOSYS}$, and \FULL{}
closes $\numslot{HSKILL-NOSYS-PCT}$ of the \NOSYS{} gap.

\begin{resultbranch}{works}
The design skill improves UI quality on both oracle signals: \FULL$\succ$\NEUTRAL{} at
$\numslot{HSKILL-POC}$ POC and $\numslot{HSKILL-WIN}$ preference, both surviving Holm. Because the
gain is measured against a token-matched, design-irrelevant prompt, it is attributable to the
skill's \emph{content} rather than to added prompt mass---the near-zero \NEUTRAL--\NOSYS{} contrast
($\numslot{NEU-NOSYS}$) confirms that filler tokens alone do essentially nothing. The rest of
Stage~1 is well-posed: there is a real effect to decompose.
\end{resultbranch}
\begin{resultbranch}{null}
\FULL$\approx$\NEUTRAL{} on both signals with a tight CI around zero: the model cannot make use of a
structured design skill at all. This outcome would have been caught at the pilot gate and escalated
to the recorded model fallback; were it to survive to the full run, we would report it honestly and
scope the remainder of Stage~1 to ``no measurable skill effect on this model,'' a negative result
about the model's instruction-following rather than about design guidance in general.
\end{resultbranch}

\subsection{RQ1 --- component necessity}
Necessity is the paired contrast \FULL{} vs.\ \LOO-$C_i$: does removing $C_i$ from the intact skill
degrade quality? Per-component objective deltas are in Table~\ref{tab:s1effects} and the twin-panel
forest plot (Figure~\ref{fig:F1}); the directional prior is that C5 (negatives) and C1 (color) punch
above their token weight.

\begin{resultbranch}{necessary (ranked)}
Removing certain components significantly degrades quality; ranked by standardized $\delta$, the
skill's weight concentrates in \numslot{NEC-RANK}, with \numslot{NEC-NSURV} of five components
surviving Holm on the POC. The largest necessity effects are C5 (negatives, $\numslot{NEC-C5}$;
BT utility $\numslot{NEC-C5-BT}$) and C1 (color, $\numslot{NEC-C1}$; BT $\numslot{NEC-C1-BT}$),
confirming the prior that the skill's leverage lies disproportionately in \emph{telling the model
what not to do} and \emph{pinning a palette}---precisely the components whose token weight is
smallest relative to their effect. The ordering replicates on held-out prompts (top component
$\numslot{NEC-REPL-TOP}$), so it is not a dev artifact. Per Eq.~\ref{eq:necsuff}, components with
large necessity but smaller sufficiency (below) carry positive interaction load---they help most in
the presence of the others.
\end{resultbranch}
\begin{resultbranch}{null}
No leave-one-out contrast degrades quality: every \FULL--\LOO-$C_i$ CI brackets zero
(Table~\ref{tab:s1effects}), and the power gate certifies these as tight nulls, not under-powering.
The reading is that no single component is individually necessary---the skill works only as a whole,
with each component substitutable by the others' combined effect. This makes RQ2/RQ3 the load-bearing
questions: a skill that is robust to any single deletion is either highly redundant or purely
interactive.
\end{resultbranch}
\begin{resultbranch}{negative}
Removing a component \emph{improves} quality: at least one \FULL--\LOO-$C_i$ contrast is significantly
\emph{negative} (\LOO{} beats \FULL). This is a genuine, honestly reported finding---an
over-constraining instruction that, in this model, does more harm than good (a plausible candidate is
an over-specified typography or pattern directive that pushes the model into awkward choices). The
practical lesson for skill design is that more guidance is not monotonically better; we report which
component and by how much, and recommend its revision.
\end{resultbranch}

\subsection{RQ2 --- component sufficiency}
Sufficiency is \AOI-$C_i$ vs.\ \NEUTRAL: does $C_i$ \emph{alone} lift quality, and does any single
component reach \FULL? The necessity$\times$sufficiency scatter (Figure~\ref{fig:F2}) reads the two
questions jointly against the $y{=}x$ line where interactions vanish. The largest single-component
lift reaches $\numslot{SUFF-MAXPCT}$ of the \FULL--\NEUTRAL{} gap.

\begin{resultbranch}{distributed}
Every component alone lifts quality above \NEUTRAL{} (\numslot{SUFF-NSURV} of five Holm-significant),
but none reaches the preregistered $30\%$ of the \FULL--\NEUTRAL{} gap---the largest,
\numslot{SUFF-TOP}, attains only $\numslot{SUFF-MAXPCT}$. Sufficiency is therefore \emph{distributed}:
no single instruction carries the skill, and the full effect emerges from the components together.
Read with RQ1 through Eq.~\ref{eq:necsuff}, components that are more necessary than sufficient (above
the $y{=}x$ line in Figure~\ref{fig:F2}) reveal positive synergy---the whole exceeds the sum of parts.
\end{resultbranch}
\begin{resultbranch}{one-carries}
A single component reaches within $30\%$ of \FULL{} on its own: \numslot{SUFF-TOP} attains
$\numslot{SUFF-MAXPCT}$ of the gap, and no other component approaches it. The skill is, to first order,
that one instruction; the remainder is polish. This is a strong, compressible result---most of a
391-token skill's benefit is reproducible with one component---and it sharpens the Stage-2 target: the
extracted direction should correspond most closely to that component's activation signature.
\end{resultbranch}
\begin{resultbranch}{null}
No component alone beats \NEUTRAL{} (all \AOI-$C_i$ CIs bracket zero) even though the intact skill
does (H-skill). Components are effective only in combination: sufficiency is strictly emergent. This
pairs with the RQ3 interaction story---if no main effect is individually sufficient yet the skill
works, the effect lives in the interactions, and the factorial block is where it is found.
\end{resultbranch}

\subsection{RQ3 --- interactions and additivity}
The 16-run factorial estimates all ten two-factor interactions cleanly (each aliased only with an
assumed-negligible three-factor term). The 2FI heatmap is Figure~\ref{fig:F3}; the directional prior
is a positive C1$\times$C5 (negatives help more once a palette is set) and a mild negative C2$\times$C3.

\begin{resultbranch}{interacting}
\numslot{INT-NSURV} two-factor interactions survive BH and are concordant on a second metric family.
The largest is \numslot{INT-MAXPAIR} ($\numslot{INT-MAX}$); as predicted, C1$\times$C5 is
super-additive ($\numslot{INT-C1C5}$)---negative constraints deliver more once a deliberate palette is
specified, because ``avoid purple gradients'' bites hardest when a replacement palette exists---while
C2$\times$C3 is mildly sub-additive ($\numslot{INT-C2C3}$). The components are not independent knobs;
the skill is a small interacting system, and this is the quantitative content of the
necessity$>$sufficiency asymmetry in Eq.~\ref{eq:necsuff}.
\end{resultbranch}
\begin{resultbranch}{additive}
No two-factor interaction survives BH: the component effects simply add, and the skill behaves as a
\emph{checklist} whose value is the sum of its independent items. This is itself informative for skill
design---components can be added, removed, or reordered without surprise---and it means necessity and
sufficiency coincide (Eq.~\ref{eq:necsuff} with $\sum_j\beta_{ij}\approx0$), which Figure~\ref{fig:F2}
would show as points on the $y{=}x$ line.
\end{resultbranch}
\begin{resultbranch}{dominated by a single main effect}
One component's main effect swamps all interactions and most other mains: the skill is effectively one
lever (\numslot{SUFF-TOP}), with the remaining components contributing little independently or
interactively. The paper's Stage-1 story then collapses to ``the skill is essentially component $X$,''
and Stage~2 asks whether $X$'s activation signature is the whole design direction.
\end{resultbranch}

\subsection{RQ4 --- do five words match 391 tokens?}
We compare the intact skill to a five-word nudge (\BEAUTY: ``Make it beautiful and well-designed.'')
and to \NEUTRAL, testing the ordered contrasts. \BEAUTY{} reaches $\numslot{NUDGE-BPCT}$ of the
\FULL--\NEUTRAL{} gap.

\begin{resultbranch}{content}
The ordering \FULL$>$\BEAUTY$>$\NEUTRAL{} holds with \FULL$-$\BEAUTY${>}0$ ($\numslot{NUDGE-FB}$ POC;
\FULL{} preferred over \BEAUTY{} at $\numslot{NUDGE-FBWIN}$): a nudge helps ($\numslot{NUDGE-BN}$ over
\NEUTRAL) but structured guidance is load-bearing and not substitutable by a one-liner. What the
391 tokens buy beyond ``make it beautiful'' is the \emph{specific} content---the palette discipline
and the negative constraints RQ1 finds most necessary---which a generic exhortation cannot supply.
\end{resultbranch}
\begin{resultbranch}{nudge}
\BEAUTY{} reaches $\ge70\%$ of the gap ($\numslot{NUDGE-BPCT}$) and \FULL$-$\BEAUTY{} is null: a
five-word nudge captures most of the skill's effect. This is a striking, cheap result---the model
already ``knows'' how to make a better page and mostly needs \emph{permission} to leave the
high-probability center, not 391 tokens of instruction. It reframes the design skill as an
activation trigger rather than a knowledge transfer, which is exactly the hypothesis Stage~2 tests
mechanistically.
\end{resultbranch}
\begin{resultbranch}{neither}
\BEAUTY$\approx$\NEUTRAL: a vague nudge does nothing, and only the specific structured skill helps (if
H-skill fired). The model needs concrete, itemized guidance or none at all---consistent with a picture
in which ``beautiful'' is too diffuse a target to move the output distribution, whereas named colors,
type scales, and prohibitions are actionable.
\end{resultbranch}

\subsection{RQ7 (black-box slice) --- does the skill effect generalize across task types?}
Before any steering, we report the \FULL--\NEUTRAL{} POC delta \emph{per task type}, including the two
held-out unseen types (settings-panel, admin-data-table). The mean seen-type delta is
$\numslot{BB-SEEN}$ and the mean unseen-type delta is $\numslot{BB-UNSEEN}$. A skill effect that holds
on the information-dense unseen types is evidence the guidance encodes transferable design content, not
a marketing-page idiom; a collapse on those types localizes the skill's benefit to the layouts it was
written for and sets expectations for the Stage-2 out-of-distribution steering test.

\begin{table}[tbp]
\centering\footnotesize
\caption{Stage-1 per-component effects on the POC (standardized $\delta$, 95\% cluster-bootstrap CI),
necessity (\FULL$-$\LOO-$C_i$) and sufficiency (\AOI-$C_i-$\NEUTRAL). Interactions and preference
utilities are reported in Figures~\ref{fig:F1}--\ref{fig:F3}. Every entry awaits the GPU run.}
\label{tab:s1effects}
\begin{tabular}{@{}lccl@{}}
\toprule
component & necessity $\delta$ (\FULL$-$\LOO) & sufficiency $\delta$ (\AOI$-$\NEUTRAL) & prior \\
\midrule
C1 color        & \numslot{NEC-C1} & \numslot{SUFF-C1} & large \\
C2 layout       & \numslot{NEC-C2} & \numslot{SUFF-C2} & medium \\
C3 typography   & \numslot{NEC-C3} & \numslot{SUFF-C3} & medium \\
C4 patterns     & \numslot{NEC-C4} & \numslot{SUFF-C4} & near-null \\
C5 negatives    & \numslot{NEC-C5} & \numslot{SUFF-C5} & largest \\
\bottomrule
\end{tabular}
\end{table}

\synthfig[0.61\linewidth]{F1}{F1_preview.pdf}{5}{Which components carry the skill? Standardized
$\delta$ vs.\ reference for each \LOO{} (necessity), \AOI{} (sufficiency), and \BEAUTY{} cell, with
Holm CIs; filled markers survive Holm. The twin-panel POC/BT layout is final; the effect sizes are
planted (C5, C1 largest; C4 near-null by construction).}

\synthfig[0.39\linewidth]{F2}{F2_preview.pdf}{5}{Necessity $\times$ sufficiency. Points above the
$y{=}x$ line are more necessary than sufficient---positive interaction load (Eq.~\ref{eq:necsuff}).
Planted coordinates; the real figure places the five components by their measured effects.}

\synthfig[0.42\linewidth]{F3}{F3_preview.pdf}{5}{Two-factor interaction coefficients (factorial
MixedLM), BH-significant cells starred. Planted C1$\times$C5 super-additivity and C2$\times$C3
sub-additivity illustrate the encoding; real signs and magnitudes replace them.}
